\section{A joint extraction at several depths}\label{sec:joint}

The constituent-stage analysis of
\cite[Theorem~6.4]{MoreAsymmetry2026} combines differently prescribed
populations at a common depth.  The extraction below also permits
populations at different depths to share one hashing step, provided
their priority orders agree on the physical tensor sides.  Population
data, strand lengths, and the constant sums defining their coarse
supports may differ.  The successive compatibility checks follow
\cite[Section~4.1]{MoreAsymmetry2026}, building on the asymmetric
hashing and interface-tensor methods of
\cite{DuanWuZhou2023,VXXZ2023}.  We specify the intended output first,
then derive the three capacities that govern its extraction yield.

For a probability vector $p$ on a finite alphabet, write
$H(p)=-\sum_a p(a)\log p(a)$, with $0\log0=0$.  A word has
\emph{type} $p$ if its empirical distribution is exactly $p$; a type window
of width $\theta$ means coordinatewise distance at most $\theta$.
All alphabets, populations, and positive population masses in this
section are fixed as $N$ tends to infinity.

\subsection{The data of a joint step}

Index the populations by a finite set $\mathcal R$.  Population $r$ has
$m_r=c_rN$ copies of a fixed strand tensor $F_r$ with finite variable
sets, where $c_r>0$ is rational.  On each side its variables are
partitioned into coarse blocks indexed by $\{0,\ldots,16\}$.
Every monomial of $F_r$ has its triple of coarse indices in an
ambient support $\Omega_r$, whose elements $u=(u_X,u_Y,u_Z)$ satisfy
\begin{equation}\label{eq:coarse-sum}
 u_X+u_Y+u_Z=s_r.
\end{equation}
Here $X,Y,Z$ is one common ordering of the physical sides.  In the
product $\bigotimes_r F_r^{\otimes m_r}$, a coarse block is specified
by a sequence of indices, one for every strand.
A fine side variable is a complete basis variable of the product input;
its coarse block records the coarse index at every strand position.
The fixed strand tensor in a population is repeated identically at
all its positions.  Populations with different histories or physical
placements are kept distinct.

Choose a rational law $p_r$ on $\Omega_r$ that maximizes entropy among
all laws on $\Omega_r$ with its three marginals.  The support of $p_r$
may be smaller than $\Omega_r$.  There are two kinds of population.

\begin{enumerate}
\item An \emph{ordinary population} delivers the full tensor block at
each selected triple $u$, without further output type tests.  The
blocks may themselves be tensor products of child strands.  This
includes extraction by whole-strand shape and the final split into
length-one strands.
\item A \emph{sharing population} consists of copies of $T_{\ell,g}$,
with even $\ell$ and $\sum_Wg_W=2\ell$.  Its coarse index $u_W$ is
the weight sum in the left half, so
\[
 \Omega_r=\{u:0\le u\le g,\ \textstyle\sum_Wu_W=\ell\},
 \qquad T_{\ell,g}|_u=T_{\ell/2,u}\otimes T_{\ell/2,g-u}.
\]
Assume $p_r(u)=p_r(g-u)$.  On each side $W$, a statistic of the full
half-strand variable takes values in a fixed finite alphabet
$\mathcal A_{r,W}$.  A map $\rho_W$ from this alphabet to weights
determines the half weight sum.  For every half-shape $u$ with
$p_r(u)>0$, prescribe a rational law $\nu_{r,u,W}$ supported on
$\rho_W^{-1}(u_W)$.

When a half-shape has a zero coordinate, the full variables on the
other two sides match by complementation.  Require that their
statistics also match by a fixed bijection and that their prescribed
laws are its pushforwards.  A side having a unique possible variable
has its corresponding point-mass prescription.  These requirements
apply as well when two coordinates vanish.
\end{enumerate}

\paragraph{Ideal outputs and inherited masks.}
A \emph{target triple} $e$ is a triple of coarse sequences having exact
joint type $p_r$ in each population.  Start with the tensor block at $e$.
In every sharing population, impose a positive type window around
$\nu_{r,u,W}$ separately in every class of positions of given
$(r,u,\text{half})$, on each side.  The right-half prescription in the
class indexed by $u$ is $\nu_{r,g-u,W}$.  Leave ordinary blocks
unrestricted.  The resulting tensor product is the \emph{ideal output}
$Q_e$.  Target triples have the same class sizes, so
their ideal outputs are isomorphic by position permutations preserving
the physical sides; denote a canonical such tensor by $Q_N$.
A fine side variable in a coarse block of $e$ is \emph{useful for $e$}
if it satisfies all these output windows; ordinary populations
contribute no further condition.  Thus usefulness specifies the
variables of the ideal tensor.  Additional input or separation tests
may delete some of them.

The child windows constrain the two halves separately.  Inherited
conditions can also constrain how their statistics are paired at parent
positions.  For a sharing population, suppress $r$ temporarily and define
the mixture law, omitting all zero-mass shapes as we do below:
\begin{equation}\label{eq:joint-pair-law}
 D_W=\sum_{u\in\Omega_r}p(u)\,
       \nu_{u,W}\otimes\nu_{g-u,W}.
\end{equation}
This is the expected pair law when, within each shape class, the left
and right statistics have their prescribed laws and are independently
arranged.

An input mask is always a deletion of individual side variables.
For sharing populations we allow finitely many empirical tests on
the ordered pair of half statistics, or a fixed function of that pair.
Their centers must be the law $D_{r,W}$ in
\eqref{eq:joint-pair-law}, or its corresponding pushforward, and their
widths must be fixed positive numbers.  Ordinary input masks must
already follow from the exact target coarse types and the specified
ideal blocks.  No other inherited masks are assumed.  All masks are
imposed separately in the populations just defined.  A mask on a
larger received class can instead be enforced separately, with
tighter widths and the same center, on any fixed subdivision; the
resulting partwise restrictions imply the original one.

Choose output widths small enough that replacing each prescribed
child law by any law in its window leaves the expected ordered-pair
law strictly inside every inherited window.  This condition has a
uniform positive margin: for coordinatewise child errors at most
$\theta$, a product-law coordinate changes by at most $2\theta$,
and all relevant collections are finite.  The same margin will be
required for the additional pair tests introduced in the proof.

\paragraph{Compatibility groups and capacities.}
The extraction will first separate targets by their $X$ coarse blocks
and impose the $X$ output windows.  It then separates $Y$ variables,
imposes their output windows, and finally separates $Z$ variables.
The later-side separation tests use only those individual shape types
that matching forces from an earlier side.  Define their designated
shapes by
\begin{equation}\label{eq:designated-compatibility}
 \mathcal I_Y=\{u\in\supp p:u_Z=0\},\qquad
 \mathcal I_Z=\{u\in\supp p:u_Xu_Y=0\}.
\end{equation}
A zero $Z$ coordinate lets matching transfer the $X$-side output
type condition to $Y$.  For $Z$, a zero
$Y$ coordinate transfers the condition from $X$; a zero $X$ coordinate
instead uses $Y$, after its output windows have been imposed.  If both
$X$ and $Y$ coordinates vanish, the $Z$ variable is forced.
For the remaining shapes with a given $W$ coordinate, the pair law
$D_W$ determines only an aggregate statistic law after the designated
contributions are removed.  We impose no individual compatibility
test on these shapes.  The exact predicate and its tolerances are
given in the incidence count below.
For $W\in\{Y,Z\}$, the partition $\mathcal G_W$ consists of the
individual groups $\{u\}$ for $u\in\mathcal I_W$, and the residual
groups $\{u\in\supp p\setminus\mathcal I_W:u_W=i\}$.
Empty groups are omitted.  A residual group is treated as residual
even when it contains just one shape.  Put $p(G)=\sum_{u\in G}p(u)$
and define
\begin{equation}\label{eq:joint-group-entropy}
 J_W=\sum_{G\in\mathcal G_W}p(G)
 H\left(\frac{\sum_{u\in G}p(u)\nu_{u,W}}{p(G)}\right),
 \qquad W\in\{Y,Z\}.
\end{equation}
The entropy $J_W$ measures the choices for one half-statistic word
within these groups.  Counting both halves gives the cost $2J_W$;
comparison with the pair-word entropy $H(D_W)$ gives the saving that
controls the number of compatible targets.  Accordingly, define the
sharing capacities by
\begin{equation}\label{eq:sharing-capacities}
 C_{r,X}=H(p_{r,X}),\qquad
 C_{r,Y}=H(D_{r,Y})-2J_{r,Y},\qquad
 C_{r,Z}=H(D_{r,Z})-2J_{r,Z}.
\end{equation}
Only the law of half-shapes is assumed invariant under $u\mapsto g-u$;
the two prescriptions $\nu_{u,W}$ and $\nu_{g-u,W}$ may differ.

For an ordinary population no compatibility restriction is needed,
and the capacities are simply
\[
 C_{r,X}=H(p_{r,X}),\qquad C_{r,Y}=H(p_{r,Y}),\qquad
 C_{r,Z}=H(p_{r,Z}).
\]
The counting arguments below give a common interpretation: each
$C_{r,W}$ is the entropy saving in the bound on competitors at side $W$.
The joint step adds these savings over populations before taking their
minimum.

\begin{theorem}[Heterogeneous joint step]\label{thm:joint-step}
Let $B_N$ be the restriction of $\bigotimes_r F_r^{\otimes m_r}$
by the inherited masks just specified, with the population data above.  Set
\begin{equation}\label{eq:joint-totals}
 A_{\rm amb}=\sum_{r\in\mathcal R}c_rH(p_r),\qquad
 C_W=\sum_{r\in\mathcal R}c_rC_{r,W},\qquad
 E=\min\{C_X,C_Y,C_Z\}.
\end{equation}
If $E>0$, then for every $\varepsilon>0$ one can choose positive
output widths, as small as desired, such that for all sufficiently
large $N$ divisible by the fixed type denominators, a direct sum of
$\exp(o(N))$ independent copies of $B_N$ has a linear restriction
to at least $\lfloor\exp(N(E-\varepsilon))\rfloor$ independent
copies of $Q_N$.

The choices of populations, rational laws, positive widths, and
divisibility requirements precede the limit $N\to\infty$.
In particular, the assertion does not require common strand lengths.
\end{theorem}

We first bound hash survival and assignment collisions.  The sequential
$X,Y,Z$ tests then put surviving monomials in a direct sum of variable
restrictions of the ideal outputs.  Section~\ref{sec:recovery} completes
the proof by bounding the missing fraction on every full-variable orbit
and restoring the ideal outputs from $\exp(o(N))$ input copies.

\subsection{Entropy and a common modular hash}

For a fixed alphabet $\mathcal A$, the number of words of type $p$ is
\begin{equation}\label{eq:type-count}
 \frac{m!}{\prod_a(mp(a))!}
   =\exp\bigl(mH(p)+O(\log(m+2))\bigr).
\end{equation}
Indeed, comparing the sum $\sum_{j=1}^n\log j$ with its integral
gives $\log(n!)=n\log n-n+O(\log(n+2))$.
There are at most $(m+1)^{|\mathcal A|}$ possible types.
Both bounds are uniform on each fixed alphabet, including types
with zero coordinates.

One convenient way to meet the entropy hypothesis is to use positive
weights $a_i,b_j,c_k$ and set
\[
 p(i,j,k)=Z^{-1}a_i b_j c_k\quad ((i,j,k)\in\Omega_r).
\]
For any competing law $q$ with the same marginals,
$-\sum q\log p=-\sum p\log p=H(p)$, because $\log p$ is a sum
of three one-coordinate functions and a constant.  Further,
\[
 \sum_{q(u)>0}q(u)\log\frac{p(u)}{q(u)}
 \le \sum_{q(u)>0}(p(u)-q(u))\le0,
\]
by $\log t\le t-1$.  Thus $H(q)\le H(p)$.
Positive rational weights give rational laws with this same property.
When the marginals specify a unique law, maximality also holds,
including boundary laws with zero probabilities.

First keep, on each side, just the coarse sequences having marginal
type $p_{r,W}$ in every population.  Call a triple satisfying these
three filters and the support conditions \emph{ambient}.  The total
number of target triples is
\begin{equation}\label{eq:target-count}
 |\mathcal T_N|=\exp\bigl(NA_{\rm amb}+O(\log N)\bigr).
\end{equation}
For a fixed $X$ coarse sequence, each possible ambient joint type
$q_r$ has the prescribed marginals.  Its number of completions in
population $r$ is
\[
 \frac{\prod_i(mp_{r,X}(i))!}{\prod_u(mq_r(u))!}
 =\exp\bigl(m(H(q_r)-H(p_{r,X}))+O(\log N)\bigr).
\]
Summing over the polynomially many types and using maximality yields
the degree bound
\begin{equation}\label{eq:ambient-x-degree}
 \#\{\text{ambient triples at a fixed $X$ block}\}
 \le \exp\bigl(N(A_{\rm amb}-C_X)+O(\log N)\bigr).
\end{equation}
Here and below the coarse counts ignore all fine masks, thereby
bounding every interaction that could actually occur.

Fix a small slack $\delta>0$.  Since $C_X\le A_{\rm amb}$ and
$E\le C_X$, choose a power $M=37^k$ with
\begin{equation}\label{eq:hash-size}
 \log M=N(A_{\rm amb}-E+\delta)+O(1).
\end{equation}
There is a set $U\subset\mathbb Z/M\mathbb Z$ with
\begin{equation}\label{eq:ap-free-size}
 \log|U|=\log M-O(\sqrt{\log M})
\end{equation}
and no nonconstant three-term arithmetic progression.  Here is a
version of Behrend's construction~\cite{Behrend1946} that also rules
out modular wraparound.  Put
$h=\lfloor\sqrt{\log M}\rfloor$, $b=\lfloor M^{1/h}/2\rfloor$,
and encode $d\in\{0,\ldots,b-1\}^h$ as
$n(d)=\sum_{j=0}^{h-1}d_j(2b)^j$.  Choose a largest class with
fixed $\sum_jd_j^2$.  Every encoded integer is less than
$(2b)^h/2\le M/2$, so a modular progression among them is an
integer progression.  In $n(d)+n(e)=2n(f)$ all digit sums are at
most $2b-2$, so there is no carry and $d+e=2f$.
Equal squared norms now give
$\|d-e\|^2=2\|d\|^2+2\|e\|^2-\|d+e\|^2=0$.
There are at most $1+h(b-1)^2$ norm classes, whence
$|U|\ge b^h/(1+hb^2)$.  The inequalities
$b\ge M^{1/h}/4$ and $b\le M^{1/h}$ imply
\eqref{eq:ap-free-size}.

Concatenate all $d=\sum_r m_r$ strand positions to form coarse words
$I,J,K$.
Let $S$ have entry $s_r$ at positions in population $r$.  Every
ambient triple satisfies $I+J+K=S$, even though the entries of $S$
can vary.  Choose independent uniform offsets $a,b'$ and a uniform
vector $v$ over $\mathbb Z/M\mathbb Z$, and set
\begin{equation}\label{eq:joint-hashes}
 x(I)=a+\langle v,I\rangle,\quad
 y(J)=b'+\langle v,J\rangle,\quad
 z(K)=\frac{a+b'+\langle v,S-K\rangle}{2}.
\end{equation}
Division by $2$ is defined since $M$ is odd.  Keep only coarse blocks
whose hashes lie in $U$.  Since $x+y=2z$, a surviving triple has
$x=y=z$.

For a fixed ambient triple $e$, let $S_e$ be its full hash-survival
event.  For every $v$ and $t\in U$ there is exactly one offset pair
giving $x=y=z=t$, so
\begin{equation}\label{eq:hash-survival}
 \Pr(S_e)=h_M:=\frac{|U|}{M^2}.
\end{equation}
Conditional on $S_e$, $(v,t)$ is uniform on
$(\mathbb Z/M\mathbb Z)^d\times U$.  A second triple with
differences $(\Delta I,\Delta J,\Delta K)$ has hashes
\[
 t+\langle v,\Delta I\rangle,\quad
 t+\langle v,\Delta J\rangle,\quad
 t-\tfrac12\langle v,\Delta K\rangle.
\]
If it shares $X$ with $e$, survival is equivalent to
$\langle v,\Delta J\rangle=0$; if it shares $Y$ or $Z$, it is
equivalent to $\langle v,\Delta I\rangle=0$.  In the shared-$Z$
case use $\Delta J=-\Delta I$.  Distinct triples cannot share two
blocks, by their common sum $S$.  Thus the tested difference vector
has a nonzero coordinate of absolute value at most $16$, which is
a unit modulo $37^k$.  Conditioning on the other coordinates of
$v$ shows, in all three cases, that
\begin{equation}\label{eq:hash-collision}
 \Pr(S_{e'}\mid S_e)=\frac1M
 \quad\text{for distinct triples sharing a block}.
\end{equation}
No further hash-dependent event is included in this conditioning.

\subsection{The finite compatibility count}

In a sharing population, prefilter side $W\in\{Y,Z\}$ by requiring
its empirical ordered-pair law to be within width $\eta$ of $D_W$.
A fine side variable is \emph{compatible} with a target triple if
its statistic type is within width $\chi$ of the prescribed law
in every designated shape class, separately on the two halves.
On the left these are the positions with $u\in\mathcal I_W$, using
$\nu_{u,W}$; on the right they are the positions with
$g-u\in\mathcal I_W$, using $\nu_{g-u,W}$.
No individual compatibility test is imposed on a residual group,
including a residual group with one shape.
The positions of each tested class are read from the target triple.
Ordinary populations impose no compatibility test.  A variable for
the whole product is compatible exactly when it is compatible in
every population.

\begin{lemma}[Compatibility incidence count]\label{lem:compatibility-count}
For fixed population data, there is a function
$\xi(\eta,\chi)\to0$ as $(\eta,\chi)\to(0,0)$ such that any
fine side variable passing the pair filters is compatible with at most
\begin{equation}\label{eq:compatibility-degree}
 \exp\bigl(N(A_{\rm amb}-C_W+\xi(\eta,\chi))
                    +O(\log N)\bigr),\qquad W\in\{Y,Z\},
\end{equation}
target triples containing its coarse block, before hashing.
The estimate is uniform in the variable and the target arrangement.
\end{lemma}

\begin{proof}
Consider first one sharing population of size $m$, and fix the
coarse $W$ word $k=(k_1,\ldots,k_m)$.  If $k$ does not have type
$p_W$, no target contains it and the bound is immediate.  Assume
henceforth that it has this type, and let $\mathcal C$ be the
set of shape words $(u_1,\ldots,u_m)$ of type $p$ with $u_{j,W}=k_j$.
Such a word determines a complete coarse triple.  If the given fine
variable has pair-statistic word $z=((a_j,b_j))_{j=1}^m$ and empirical
law $D'$, let $\mathcal B$ be its orbit under
\[
 \Gamma_k=\prod_i\operatorname{Sym}\{j:k_j=i\}.
\]
This group acts transitively on $\mathcal B$ by definition and on
$\mathcal C$ because all shape counts within every coarse class are
fixed.  Compatibility is unchanged by applying the same permutation
to a candidate and a statistic word.  Hence its incidence degrees
$a=\#\{z\in\mathcal B:z\text{ compatible with }c\}$ and
$b=\#\{c\in\mathcal C:z\text{ compatible with }c\}$ do not depend
on $c$ and $z$, respectively.  Counting incidences gives the exact
finite identity
\begin{equation}\label{eq:incidence-identity}
 |\mathcal C|a=|\mathcal B|b.
\end{equation}
Since the statistic determines the half weight,
\begin{align*}
 |\mathcal C|&=\frac{\prod_i(mp_W(i))!}{\prod_u(mp(u))!},
 &\log|\mathcal C|&=m(H(p)-H(p_W))+O(\log(m+2)),\\
 |\mathcal B|&=\frac{\prod_i(mp_W(i))!}{\prod_{a,b}(mD'(a,b))!},
 &\log|\mathcal B|&=m(H(D')-H(p_W))+O(\log(m+2)).
\end{align*}

To bound $a$, now fix a candidate $c\in\mathcal C$.  This fixes the
position set of every left and right compatibility group.  On the
left, compatibility specifies the law within each designated group
to accuracy $\chi$.  The row sums of $D'$ give the total count of
each statistic among positions with each coarse value, since
$\rho_W(a_j)=k_j$.  Subtracting the designated-group counts leaves, in the
remaining group with coarse value $i$, the law
\[
 \frac{\sum_{u\in G}p(u)\nu_{u,W}}{p(G)}
\]
up to an error tending uniformly to zero with $\eta$ and $\chi$.
More explicitly, each coordinate error is at most
$(|\mathcal A_{r,W}|\eta+\chi)/p(G)$: summing a row of $D'$
costs at most $|\mathcal A_{r,W}|\eta$, and the weighted errors
of the subtracted designated laws total at most $\chi$.
Only fixed positive numbers $p(G)$ are divided by.  A product of
multinomial counts over these fixed position sets, and then a sum
over polynomially many types, bounds the number of left-statistic
assignments by $\exp(m(J_W+o_{\eta,\chi}(1))+O(\log(m+2)))$.

For the right half, use the column sums of $D'$.  Its coarse value
is $g_W-k_j$, so those sums determine precisely the analogous
aggregate counts.  The right half-shape distribution is
$p(g-u)=p(u)$, giving the same entropy $J_W$, with the actual right
prescription $\nu_{g-u,W}$.  Count the two halves independently;
discarding the requirement that they reassemble the exact pair law
$D'$ can only enlarge the count.  We obtain
\[
 a\le\exp\bigl(m(2J_W+o_{\eta,\chi}(1))+O(\log(m+2))\bigr).
\]
There is no factor for choosing group locations: those locations
were fixed with $c$.  By \eqref{eq:incidence-identity} and uniform
continuity of entropy on a finite simplex,
\[
 b\le\exp\bigl(m(H(p)-H(D_W)+2J_W+o_{\eta,\chi}(1))
                         +O(\log(m+2))\bigr).
\]
For an ordinary population all candidates are compatible, and their
number is $\exp(m(H(p)-H(p_W))+O(\log(m+2)))$ directly.  The
candidate sets, permutation groups, and compatibility conditions
factor over populations.  Multiplying their bounds proves
\eqref{eq:compatibility-degree}.
\end{proof}

\subsection{Sequential assignment by variable deletion}

Choose $\eta,\chi>0$ sufficiently small that the error in
Lemma~\ref{lem:compatibility-count} is below $\delta/4$.  Next
choose output widths $\theta>0$ smaller than $\chi$ and small
enough for the positive margins required by all inherited and pair
tests.  Complementary prescribed laws ensure that usefulness of
one side transfers to the compatibility window on a matching side.
Every choice here is fixed before $N$ grows.

First delete an $X$ coarse block unless it belongs to exactly one
ambient hash survivor and that survivor is a target triple.  Call
the retained triples \emph{eligible}.  Freeze this list for all
subsequent tests.  By \eqref{eq:ambient-x-degree},
\eqref{eq:hash-size}, and \eqref{eq:hash-collision}, a target triple
$e$ fails this test, conditional only on $S_e$, with probability at
most $\exp(-\delta N+O(\log N))$.  Each retained $X$ block now has
an unambiguous assignment.  In that block retain only the fine
variables useful for its assigned triple, that is, those satisfying
its output windows.

On $Y$, apply its pair filters and delete every fine variable
having other than exactly one compatible eligible triple containing
its coarse block.  Assign it to that triple and impose usefulness
for that assignment.  Apply the same procedure to $Z$, using
the designated classes $\mathcal I_Z$.  Include all inherited masks
as side-variable deletions; their application does not change the frozen eligible
list or the definitions of the compatible lists.

We verify that these operations produce a direct sum, rather than
merely a list of proposed components.  Take an actual monomial
surviving every deletion, and let $e$ be its ambient coarse triple.
Its $X$ block survived, so $e$ is that block's unique eligible
triple and its $X$ variable is useful for $e$.  In every sharing
half with $u_Z=0$, the $Y$ statistics are complementary to the
$X$ statistics.  Thus the actual $Y$ variable is compatible with
$e$; deterministic sides satisfy their point-mass tests as well.
The $Y$ uniqueness test therefore assigns this variable to $e$.
Its usefulness is now known for the actual triple.

For $Z$, a half with $u_Y=0$ transfers the required statistics from
$X$.  A half with $u_X=0$ and $u_Y>0$ transfers them from the
already useful $Y$ variable.  With both $u_X=u_Y=0$, the remaining
variable is forced.  These cases prove compatibility with $e$ on
every designated shape class in $\mathcal I_Z$.  Its uniqueness test then
assigns the actual $Z$ variable to $e$ as well.  Ordinary
populations add no compatibility conditions in either argument.
All three variables of every retained monomial therefore have the
same assignment.

Different assignments share no retained fine variable on any side.
Within assignment $e$, usefulness restricts precisely to $Q_e$;
all other tests delete individual variables of $Q_e$.  Consequently
the surviving tensor is a direct sum of tensors $Q_e^{\rm hole}$,
each obtained from its ideal $Q_e$ solely by variable deletion.
This statement concerns complete variables of the population
product; sharing a local coordinate factor between two complete
variables does not identify those variables.

Here is the probability estimate needed to repair these holes.
For a fixed $e$ and a fixed variable $v$ of $Q_e$ passing all its
deterministic inherited and pair masks, own-triple compatibility
and usefulness hold.  Let $R_{e,v}$ be the union of the $X$
eligibility-failure event and, when $v$ is on side $Y$ or $Z$,
the events that another compatible target triple sharing its
coarse block survives the hash.  Every random deletion of $v$ is
contained in $R_{e,v}$.  Bound those competitors among
\emph{all} target triples before eligibility and other deletions.
Lemma~\ref{lem:compatibility-count} and
\eqref{eq:hash-collision} then give, for some constant $c>0$,
\begin{equation}\label{eq:joint-conditional-loss}
 \Pr(R_{e,v}\mid S_e)
 \le \exp(-cN),
\end{equation}
uniformly in $e$ and the passing variable, for all large $N$.
Indeed, the exponent for side $W$ is at most
$N(E-C_W-\delta+\xi)+O(\log N)$ and $C_W\ge E$.
This estimate is conditioned only on $S_e$, never on eligibility
or on passing a hash-dependent test.

For reference in the recovery argument, let $G_e$ act by independent
position permutations in every ideal child class $(r,u,\text{half})$,
simultaneously on all three sides.  In ordinary whole-strand blocks
use their position permutations; for a factorized ordinary block
one may also permute its child factors independently.  These are
automorphisms of $Q_e$.  The full variable alphabets in its finitely
many classes are fixed, so $G_e$ has only polynomially many orbits
of variables.  Lemma~\ref{lem:orbit-masks} supplies the uniform
deterministic mask bound on these orbits.  The completion of
Theorem~\ref{thm:joint-step} in Section~\ref{sec:recovery} combines
that bound with \eqref{eq:joint-conditional-loss}, selects
exponentially many good assignments using
\eqref{eq:target-count} and \eqref{eq:hash-survival}, and applies
Lemma~\ref{lem:hole-repair} to reconstruct their ideal outputs.
